\documentclass[a4paper]{article}
\input{preamble.tex}
\title{Analysis 2: Midterm 1}
% \author{Aamod Varma}
\begin{document}
\maketitle
\textbf{Problem 1.}

\vspace{1em}
Let $$g_i(x) = \begin{cases}
    \sin(x) \quad &2\pi (i - 1) \le x < 2\pi i\\
    0 \quad & \text{ otherwise }
\end{cases}$$

Now define,
\[
    f(x) = \sum_{i = 1}^{\infty} \frac{1}{i} g_i(x)
\]


First, clearly $f$ is bounded as for each interval of $2\pi$ starting from $0$ as it's defined as $g_i$ for each interval and trivially we know that $|\sin| \le 1$. Now we also know that $g_i$ itself is continuous because $\sin$ is continuous. Now we need to show that $f$ is continuous at the points $2\pi i$ where $i = 0, 1, 2, \dots$. Note that $g_i$ is defined from $[2\pi (i - 1), 2\pi i)$ so we need to test continuity at $2\pi i$. Note limit as $g_i$ approaches $2 \pi i$ is simply $0$ and similarly we have $g_{i+1}$ at $2\pi i$ as $0$. Hence, our functions is continuous.

\vspace{1em}

As $f$ is continuous and bounded it is Riemann-inntegrable on each $[0, n]$. Now we compute $\int_0^{n} f$. Note that in the range $[0, n]$ $f$ can be defined as a finite sum i.e $\sum_{i = 1}^{\lceil\frac{n}{2\pi} \rceil}\frac{1}{i}g_i(x)$ and hence we can write, 
\begin{align*}
    \int_{0}^{n} f(x) &= \int_{0}^{n} \sum_{i = 1}^{\lceil \frac{n}{2\pi} \rceil} \frac{1}{i} g_i(x)\\
                      &= \sum_{i = 1}^{\lceil \frac{n}{2\pi} \rceil} \frac{1}{i} \int_{0}^{n} g_i(x)\\
\end{align*}

Now clearly integrating each $g_i$ from $[0, n]$ is equivalent to integrating it where it's define non-zero in this domain. So we can rewrite this as,
\begin{align*}
    \int_{0}^{n} f(x) &= \sum_{i = 1}^{\lceil\frac{n}{2\pi}\rceil} \frac{1}{i} \int_{0}^{n} g_i(x)\\
                      &= \int_{0}^{2\pi} g_1(x) + \frac{1}{2} \int_{2\pi}^{4\pi} g_2(x) + \dots + \frac{1}{\lceil\frac{n}{2\pi}\rceil} \int_{2\pi (\lceil\frac{n}{2\pi}\rceil - 1)}^{n}  g_{\lceil\frac{n}{2\pi}\rceil}(x)
\end{align*}

Note that each $g_i$ integrated where it's defined is simply just zero as the positive and negative cancel each other out.  So we have, 
\[
\int_{0}^{n} f(x) = \frac{1}{\lceil\frac{n}{2\pi}\rceil} \int_{2\pi (\lceil\frac{n}{2\pi}\rceil - 1)}^{n}  g_{\lceil\frac{n}{2\pi}\rceil}(x)
\]

However, now note that as $\sin$ is bounded above by 1, we have $ \left |\int_{2\pi (\lceil\frac{n}{2\pi}\rceil - 1)}^{n}  g_{\lceil\frac{n}{2\pi}\rceil}(x)\right| \le 2$ and hence, 
\begin{align*}
    \int_{0}^{n} f(x) &= \frac{1}{\lceil\frac{n}{2\pi}\rceil} \int_{2\pi (\lceil\frac{n}{2\pi}\rceil - 1)}^{n}  g_{\lceil\frac{n}{2\pi}\rceil}(x)\\
                      &\le \frac{1}{\lceil\frac{n}{2\pi}\rceil} 2
\end{align*}

So, 
\[
    \lim_{n \to \infty} \int_{0}^{n} f(x) \le \lim_{n \to \infty}   \frac{2}{\lceil \frac{n}{2\pi}\rceil} = 0
\]

So we have the integral converges to 0. Now if we consider $|f|$ instead. And note that this is equivalent to $\sum_{n = 1}^{\infty} \frac{1}{i}|g_i(x)|$. Following the same logic we get, 
\begin{align*}
    \int_{0}^{n} |f(x)| &= \int_{0}^{n} \sum_{n = 1}^{\lceil n / 2\pi \rceil} \frac{1}{i}|g_i(x)|\\
                        &=  \sum_{i = 1}^{\lceil n / 2\pi \rceil} \frac{1}{i} \int_{0}^{n}|g_i(x)|\\
\end{align*}

However, note here that we have $\int_{0}^{n} |g_i(x)| = \int_{2\pi (i - 1)}^{2\pi i} |\sin(x)|$, $\sin$ is positive in the first half of the interval so this is equivalent to $2 \int_{2\pi (i - 1)}^{\pi} \sin(x) = 4$. So we have,

\begin{align*}
    \int_{0}^{n} |f| &\ge \sum_{i = 1}^{\lceil n / 2 \pi  \rceil -1} \frac{4}{i}  
\end{align*}

Note that we can ignore the last (possibly incomplete integral for simplicity) this gives us. But this means that, 
\[
    \lim_{n \to \infty} \int_{0}^{n} |f| &\ge 4 \lim_{n \to \infty} \sum_{i = 1}^{\lceil n / 2 \pi  \rceil -1} \frac{1}{i}
\]

Note the right hand side is the harmonic series as $n \to \infty$  and we know that it diverges.  Since we just showed that the left side is greater than the sum of the harmonic series, this means that the left side diverges. 


\vspace{1em}
\textbf{Problem 2.}

\quad Given that $f$ is $C^{1}$ on $U \subset R^{n}$. Now note that $f$ is a vector valued function so we can write it component wise as $f = (f_{1}, \dots, f_n)$ where $f_i:R^{n} \to R$. Now consider the parametrization of the line joining $x, y$ as $g(t) = y + (x - y)t$ where $0 \le t \le 1$. Now note again that in order to show, 
\[
    f(x) - f(y) = \int_{0}^{1}  f'(y + t(x - y)) \cdot (x - y) dt
\]

It is enough to show this for each component (as vector integration is defined component wise). This means we just need to show,

\[
    f_i(x) - f_i(y) = \int_{0}^{1}  f_i'(y + t(x - y)) \cdot (x - y) dt
\]

Now consider the function $f_i \circ g$ and as $g: R \to R^{n}$ and $f_i: R^{n} \to R$ we have $f_i \circ g: R \to R$ a 1D function. We can now apply the 1D fundamental theorem of calculus to this function to get, 
\begin{align*}
    (f_i \circ g)(1) -  (f_i \circ g)(0) &= \int_{0}^{1} (f_i \circ g)' \ dt\\
                f_i(g(1)) - f_i(g(0)) &= \int_{0}^{1} f_i'(g(t)) g'(t) \ dt\\
                f_i(x) - f_i(y) &= \int_{0}^{1} f_i'(y + t(x - y)) (x - y) \ dt
\end{align*}

As this is true for any component for $f$ we have, 
\begin{align*}
    f(x) - f(y) &= (f_0(x) - f_0(y), \dots, f_n(x) - f_n(y))\\
                &= \left(\int_{0}^{1} f_0'(y + t(x - y)) (x - y) \ dt, \dots \int_{0}^{1} f_1'(y + t(x - y)) (x - y) \ dt \right)\\
                &= \int_{0}^{1} f'(y + t(x - y)) \cdot (x - y) \ dt
\end{align*}        



\vspace{1em}
\textbf{Problem 3.}

\quad 1. We have $ || f'(x) - I| | \le \frac{\epsilon}{2}$, now using definition of operator norm we get, 
\[
   \frac{ \left | f'(x)v - v \right | }{|v|}\le \left | \left | f'(x) - I \right |\right | \le \frac{\epsilon}{2}
\]

Now using the theorem from Problem 2 taking $x = v$ and $y = 0$ we get, 
\begin{align*}
    f(v) - f(0) &= \int_{0}^{1} f'(tv) v \ dt\\
        f(v) &= \int_{0}^{1} f'(tv) v \ dt \\
        f(v) - v &= \int_{0}^{1} f'(tv) v - v \ dt \\
\end{align*}

Now taking the norm from both sides we have, 
\begin{align*}
    f(v) - v &= \int_{0}^{1} f'(tv) v - v \ dt \\
    |f(v) - v| &= \left | \int_{0}^{1}  f'(tv) v - v \ dt\right |\\
               &\le  \int_{0}^{1} | f'(tv) v - v |\ dt\\
\end{align*}

But note from above we have, 
\begin{align*}
    \frac{|f'(x)v - v|}{|v|} &\le \frac{\epsilon}{2}\\
    |f'(x) v - v| &\le |v| \frac{\epsilon}{2}
\end{align*}

Now as $v \in U$ we have $tv \in U$ for $t \in [0, 1]$ hence we also get, $ |f'(tv) v - v| &\le |v| \frac{\epsilon}{2}$ and that,

\begin{align*}
    |f(v) - v| \le \int_{0}^{1} |f'(tv)v - v| dt  &\le \int_{0}^{1} |v| \frac{\epsilon}{2} \ dt\\
                                                  &= \frac{|v|\epsilon}{2} [t]_0^{1} = \frac{|v| \epsilon}{2}
\end{align*}

Dividing by $|v|$ we have, 
\[
    \frac{|f(v) - v|}{|v|} \le \frac{\epsilon}{2} < \epsilon
\]

Now using this I use the triangle inequality to get $\frac{|f(v) - v|}{|v|} \ge | |f(v)| - |v| | / |v|$ so,
\begin{align*}
    \frac{| | f(v) |- |v| |}{|v|} &\le \epsilon\\
{| | f(v) | / |v|- |v|/ |v| |} &\le \epsilon\\
\left| \frac{|f(v)|}{|v|} - 1 \right| &\le \epsilon
\end{align*}

which is equivalent to $\frac{|f(v)|}{|v|} \in (1 - \epsilon, 1 + \epsilon)$



\vspace{1em}

\quad 2. We need to show that, 
\[
    B(0, 1 - \epsilon) \subset f(U) \subset B(0, 1 + \epsilon)
\]

First we need to show that for any $y \in B(0, 1 - \epsilon)$ we have some $x \in U$ such that $f(x) = y$. First consider $f(y)$ and we have, 
\begin{align*}
    |f(y) - y| < \epsilon |y|
\end{align*}


Now let we need to take a step in order to get a better estimate which is $x_{1} = y + (y - f(y))$ and now note that, 
\[
    |x_{1} - y| = |f(y) - y| < \epsilon |y|
\]

now define our steps in terms of our previous estimates such that we have $x_{n + 1} = x_n + (y - f(x_n))$. First we need to show that $f(x_k)$ is defined i.e $x_k \in U$. Note we have, 
\begin{align*}
    |x_{n} - y| &= |f(x_{n - 1}) - x_{n - 1}| \le \epsilon |x_{n - 1}|\\
    |x_{n}| &\le \epsilon |x_{n - 1}| + |y|\\
                &\le \epsilon (\epsilon |x_{n - 2}| + |y|) + |y|\\
                &\le |y| + \epsilon |y| + \epsilon^2 |y| + \dots + \epsilon^{n} |y|\\
                &= |y| (1 + \epsilon + \dots + \epsilon^{n})\\
                &= |y| \frac{1}{1 - \epsilon} \le \frac{1 - \epsilon}{1 - \epsilon} = 1
\end{align*}

So we have $|x_n| < 1$ for all $n$ which means $x_n \in U$.


\vspace{1em}

Now we need to show that with each new estimate  we get to a better solution and thus get to some limit of a solution which is the perfect solution. I.e we need to show the sequence $y, x_{1}, \dots$ has a limit. We first have, 
\begin{align*}
    |x_{n + 1} - x_n| &= |x_n + y - f(x_n) - (x_{n - 1} + y - f(x_{n - 1}))|\\
                      &= |(f(x_{n - 1}) - x_{n - 1}) - (f(x_n) - x_n))|
\end{align*}

First note that if we consider $f(x) - x$ then using FTC on this function we have, 
\begin{align*}
    (f(a) - a) - (f(b) - b) &= \int_{0}^{1}  (f(x) - x)^{(1)}(b + t(a - b)) (a - b) \ dt\\
    |(f(a) - a) - (f(b) - b)| &\le \int_{0}^{1}  \frac{\epsilon}{2} (a - b) = \frac{\epsilon}{2} (a - b)
\end{align*}


So we have, 
\[
    |x_{n + 1} - x_n|  \le \frac{\epsilon}{2} |x_n - x_{n - 1}|
\]

Using this we get, 
\[
    |x_m - x_n| < \epsilon
\]

Or that it's a Cauchy sequence and hence has a limit. But if it has a limit say $x$ then applying that to $x_{n + 1} = x_n + (y - f(x_n))$ we get $x = x + (y - f(x))$ or that $f(x) = y$ and as each $x_n \in U$ we have the limit in $U$ as well. Now this proves that $B(0, 1 - \epsilon) \subset f(U)$. We need to now show that $f(U) \subset B(0, 1 + \epsilon)$ or that for any $x$ we have $|f(x)| \le 1 + \epsilon$. 

\vspace{1em}

First note that for $x \in U$ we have $|x| \le 1$. So using 1. we get, 
\begin{align*}
    \frac{|f(x)|}{|x|} &< 1 + \epsilon\\
    |f(x)| &< (1 + \epsilon)|x| \le  1 + \epsilon
\end{align*}

So we get $|f(x)| \le 1 + \epsilon$ for all $x \in U$. 

\end{document}


